\section{Results}
\label{sec:results}

\subsection{Two development screens lead to two negative decisions}

Both equation-discovery candidates pass their development screens. Round 1
has a median development improvement of \RouteOneDevelopment{}, and Round 2
has \RouteTwoDevelopment{}. Their frozen unseen panels give a different
decision. Round 1 has a system-level 95\% interval of
[\RouteOneCILower{}, \RouteOneCIUpper{}], and Round 2 has
[\RouteTwoCILower{}, \RouteTwoCIUpper{}]. Each interval crosses zero, so both
rounds are negative under their registered criteria.

Figure~\ref{fig:route_a_results} shows these intervals before the detailed
values in Table~\ref{tab:route_a_results}.

\begin{figure}[t]
  \centering
  \includegraphics[width=\columnwidth]{figures/route-a-results.pdf}
  \Description{Two horizontal confidence intervals cross a vertical zero
  boundary, so neither equation-discovery round passes its unseen criterion.}
  \caption{Route A unseen results. Points show median failure-aware relative
  improvement and bars show frozen system-level bootstrap intervals. The
  dashed line is the acceptance boundary.}
  \label{fig:route_a_results}
\end{figure}

\begin{table}[t]
\centering
\caption{Development values and unseen decisions for the two parent-linked
mechanism rounds. Full identifiers are retained in the artifact.}
\label{tab:route_a_results}
\resizebox{\columnwidth}{!}{\input{tables/route-a-results}}
\end{table}

The tested smoothing and support-selection mechanisms therefore do not
establish improvement over Operon on the held-out systems. The state machine
keeps each favorable development value and unfavorable unseen decision,
connects Round 2 to the Round 1 result, and activates the preregistered systems
route after the second failure. This is negative scientific evidence and a
workflow transition, not a positive equation-discovery result.

\subsection{Task-level controller comparison}

Figure~\ref{fig:main_results} reports descriptive success proportions for the
three controllers. The one-shot controller succeeds on \OneShotSuccess{} of
the deterministic cells, execute-once on \ExecuteOnceSuccess{}, and the full
loop on \FullLoopSuccess{}. These values are properties of the frozen matrix.

\begin{figure}[t]
  \centering
  \includegraphics[width=\columnwidth]{figures/main-results.pdf}
  \Description{Bar chart with descriptive success proportions 0.20 for
  one-shot, 0.50 for execute-once, and 1.00 for the full loop.}
  \caption{Descriptive controller success in the controlled matrix. Each bar
  aggregates ten tasks repeated under three deterministic idempotency
  identifiers; the identifiers are not independent samples.}
  \label{fig:main_results}
\end{figure}

At the independent task level, the full-loop minus execute-once difference
vector is \([0,0,0,1,0,0,1,1,1,1]\). The mean is \TaskLevelGain{} and the
\BootstrapResamples{}-resample task bootstrap interval is
[\TaskLevelCILower{}, \TaskLevelCIUpper{}]. There are \NonTieTaskCount{}
positive differences, no negative differences, and five ties. The exact
one-sided sign-test \(p\)-value is \SignTestOneSided{}, while the two-sided
value is \SignTestTwoSided{}. The two-sided value and the wide interval do not
support a publication-facing statement that the controller performs better
on a population of research tasks. The unchanged mean is reported as an
observed contrast only.

The earlier interval [\HistoricalCellCILower{},
\HistoricalCellCIUpper{}] treated 30 deterministic task-repetition cells as
the resampling pool. It remains in the immutable parent object as evidence of
the original internal decision. It is excluded from the current abstract,
contribution statement, and scientific interpretation.

\subsection{Controlled recovery and component behavior}

The full loop begins with 24 negative cells and repairs 15, for a descriptive
recovery proportion of \FullLoopRecovery{}. Execute-once has no post-execution
action and records zero recovery. The full-loop scientific and reproduction
hashes agree for every cell, giving \FullLoopReproduction{} exact replay, and
its reports contain \FullLoopUnsupported{} unsupported claims under the
frozen checker. The campaign records \HumanInterventions{} human research
decisions after start.

Figure~\ref{fig:component_results} and Table~\ref{tab:all_modes} report the
registered component removals. Removing \vault{} memory gives
\NoVaultSuccess{} descriptive success because memory-dependent repairs remain
blocked. Removing failure feedback gives \NoFeedbackSuccess{}. Removing the
freeze gives \NoPreregSuccess{} under an endpoint that requires a frozen
contract. Removing the evidence gate gives \NoGateSuccess{} apparent success
and permits \NoGateUnsupported{} unsupported claims. These contrasts show
that the co-designed evaluator responds to each removed contract; they are
not estimates for independently authored research tasks.

\begin{figure}[t]
  \centering
  \includegraphics[width=\columnwidth]{figures/ablation-results.pdf}
  \Description{Bar chart comparing full-loop success of 1.00 with four
  component removals: no Vault 0.70, no feedback 0.50, no freeze 0.00, and no
  evidence gate 0.80.}
  \caption{Descriptive success for the complete loop and four component
  removals. The no-evidence-gate condition also emits six unsupported claims.}
  \label{fig:component_results}
\end{figure}

\begin{table*}[t]
\centering
\caption{All controller modes. Recovery is conditional on an initially failed
deterministic cell. Reproduction compares primary and rerun scientific hashes.
Unsupported is a count over 30 cells per mode.}
\label{tab:all_modes}
\small
\input{tables/mode-results}
\end{table*}

\subsection{Execution and artifact checks}

The main controller summaries create \LocalRequests{} local-model requests.
\LocalFallbacks{} request reaches its timeout and uses the declared fallback.
The event is present in policy evidence and does not affect numerical
adjudication. Paid external model and cloud-compute cost is zero.

Each matrix cell has a result JSON file, a reproduction record, a Markdown
report, and hashes for source and policy evidence. The manuscript rewrite adds
a citation-status registry, a ledger for all 28 affected publication
surfaces, full language and LaTeX scans, a compiled PDF, and a page-by-page
visual inspection record. These checks establish artifact consistency only.
